\documentclass[11pt]{article}
\usepackage[margin=1in]{geometry}
\usepackage{amsmath,amssymb,amsthm,mathtools}
\usepackage[hidelinks,pdfusetitle]{hyperref}
\usepackage{microtype}
\usepackage{array}
\newcolumntype{L}[1]{>{\raggedright\arraybackslash}p{#1}}
\newcommand{\DOI}[1]{\mbox{\href{https://doi.org/#1}{\nolinkurl{#1}}}}
\theoremstyle{plain}
\newtheorem{theorem}{Theorem}
\newtheorem{lemma}{Lemma}
\newtheorem{proposition}{Proposition}
\newtheorem{corollary}{Corollary}
\newtheorem{remark}{Remark}
\theoremstyle{definition}
\newtheorem{definition}{Definition}
\DeclareMathOperator{\Tr}{Tr}
\DeclareMathOperator{\rank}{rank}
\newcommand{\Nodd}{N_{\mathrm{odd}}}
\newcommand{\NoddE}{N_{\mathrm{odd},e}}
\newcommand{\Dop}{\Delta}
\newcommand{\dist}{\operatorname{dist}}
\setlength{\emergencystretch}{2em}
\title{A strong-coupling ceiling of $6000\,t^{2}/U$ for electronic pair
correlations in the two-dimensional Holstein model}
\author{Davis Rens\\SuperMatics Inc.\\\href{https://orcid.org/0000-0001-5839-5197}{ORCID 0000-0001-5839-5197}}
\date{}
\begin{document}
\maketitle

\begin{abstract}
We prove exponential clustering of equal-time on-site electronic pair correlations
for the spinful Holstein model on the square torus $(\mathbb{Z}/N\mathbb{Z})^2$ with
side $N\ge4$. With $U=F^{2}/K$ the only hypothesis is $\beta t^{2}/U\le 1/6000$,
that is $k_{\mathrm{B}}T\ge 6000\,t^{2}/U$; on tori of even side the constant drops
to $4500$. No restriction on the ratio $t/U$ is imposed, and the bound is uniform in
the phonon frequency, in the torus side and in the chemical potential. At distinct
sites
\[
 \bigl|\langle\Dop_x^\dagger\Dop_y\rangle\bigr|
 \le \frac{57}{500}\left(\frac{99}{100}\right)^{d(x,y)},
\]
with $\Dop_x=c_{x\downarrow}c_{x\uparrow}$ and $d(x,y)$ the torus distance, and the
rows of the correlation matrix are uniformly summable; hence the clustering
temperature $T_{\mathrm{pair}}$ of Definition~\ref{def:tpair} obeys
$k_{\mathrm{B}}T_{\mathrm{pair}}\le6000\,t^{2}/U$. The constant improves the
previously released value $2^{27}=134\,217\,728$ for the same model by a factor of
$22\,369$, and the dimension-tracking two-dimensional constant $15\,925\,248$ of the
companion preprint by a factor of $2654$, at the price of a much slower decay ratio:
$99/100$ here against $5/64$ in the earlier release and $5/24$ in the companion
preprint.
The mechanism is exact at every step: the sixteen-state single-bond spectrum with a
half-filling supremum, a bridge-parity doubling, an explicit four- or five-edge
colouring with a Schatten H\"older decoupling that costs a factor $c\le5$, exact
enumeration of connected supports with at most six edges, a four-ary Raney tree tail,
and the Koteck\'y--Preiss vertex criterion at $B=557/500$. A single certificate
program recomputes every finite witness in exact rational and symbolic arithmetic;
it also certifies that the same scheme fails at $C=5800$ on all tori and at
$C=4300$ on even tori. The bound is on a clustering temperature of one two-point
function. It is not a critical temperature, it is not a statement about any
material, and no claim of superconductivity is made.
\end{abstract}

\section{Model and main theorem}\label{sec:model}

\subsection{The spinful Holstein model on a square torus}

Fix an integer side $N\ge4$ and let $\Lambda=T_N=(\mathbb{Z}/N\mathbb{Z})^{2}$ be the
square torus with its nearest-neighbour edge set $E(\Lambda)=\{\{x,y\}\subset\Lambda:
|x-y|_{1}=1\}$ and $|E(\Lambda)|=2N^{2}$ edges, where
$|x-y|_{1}=\sum_{j=1}^{2}\min(|x_{j}-y_{j}|,N-|x_{j}-y_{j}|)$ is the $\ell_{1}$
torus distance, written $d(x,y)$ below. Each site carries a single oscillator
coordinate $X_x$ and a spinful electron level with four states, empty, up, down and
doubly occupied. The Holstein Hamiltonian \cite{holstein1959} is
\begin{equation}\label{eq:H}
H_\Lambda(\mu)=-t\sum_{e\in E(\Lambda)}B_e
+\sum_{x\in\Lambda}\left[\frac{P_x^{2}}{2m}+\frac{K}{2}X_x^{2}
-F X_x\,(n_x-1)\right]-\mu\sum_{x\in\Lambda}(n_x-1),
\end{equation}
with $B_e=\sum_{\sigma\in\{\uparrow,\downarrow\}}
(c_{x\sigma}^{\dagger}c_{y\sigma}+c_{y\sigma}^{\dagger}c_{x\sigma})$ the
spin-summed hopping across $e=\{x,y\}$, $n_x=n_{x\uparrow}+n_{x\downarrow}$, and real
parameters $t\ge0$, $m>0$, $K>0$, $F$, $\mu\in\mathbb{R}$ with
$[P_x,X_{x'}]=-i\hbar\delta_{xx'}$. Following the released analysis of the same
model we write $U=F^{2}/K>0$ for the polaronic energy,
$\Omega=\hbar\sqrt{K/m}>0$ for the phonon frequency, $\tau=k_{\mathrm{B}}T$ for the
thermal energy and $\beta=1/\tau$; there is no direct electron-electron
interaction. The chemical potential is shifted, $\mu_c=\mu-U/2$, only so that the
atomic electronic Hamiltonian in the odd-occupation gauge reads
\begin{equation}\label{eq:D0}
D_0=\frac{U}{2}\Nodd-\mu_c N,\qquad
\Nodd=\sum_{x\in\Lambda}\mathbf{1}_{\{n_x=1\}},\qquad
N=\sum_{x\in\Lambda}n_x .
\end{equation}
The physical equal-time on-site pair field is
$\Dop_x=c_{x\downarrow}c_{x\uparrow}$, where the order of the two operators is fixed
once and for all. For $\mu$, $t$, $N$, $\beta$, the grand-canonical expectation is
$\langle\,\cdot\,\rangle=\langle\,\cdot\,\rangle_{\beta,\mu,\Lambda}$; throughout,
$\tau$ is varied at fixed $\mu$, so the grand-canonical ensemble is the natural one
here and the canonical statement is added separately in
Section~\ref{sec:canonical}.

\begin{definition}[Clustering temperature]\label{def:tpair}
Let $T_{\mathrm{pair}}$ be the infimum of the temperatures $\tau$ with the following
property. For every $\tau'>\tau$, every $N\ge4$, every real $\mu$ and every
$\Omega>0$, there are constants $A_0>0$ and $\xi>0$, independent of $N$, of $\mu$ and
of $\Omega$, such that
$\bigl|\langle\Dop_x^\dagger\Dop_y\rangle_{\beta',\mu,\Lambda}\bigr|
\le A_0e^{-d(x,y)/\xi}$ for all distinct $x,y\in\Lambda$, and such that the rows of
the correlation matrix are summable uniformly in $N$,
$\sup_{N\ge4}\max_{x\in\Lambda}\sum_{y\in\Lambda}
\bigl|\langle\Dop_x^\dagger\Dop_y\rangle\bigr|<\infty$.
\end{definition}

$T_{\mathrm{pair}}$ is a clustering temperature. It is an upper threshold for the
exponential decay of one two-point function and for the summability of its rows, and
it says nothing about the existence of a phase transition at a critical temperature.
The theorem below bounds $T_{\mathrm{pair}}$ from above.

\subsection{Main theorem}\label{sec:main}

\begin{theorem}[Uniform-frequency pair-clustering ceiling in two
dimensions]\label{thm:main}
Fix $U=F^{2}/K>0$ and $t\ge0$. There is no hypothesis on the ratio $t/U$, and the
statements below are uniform in $\Omega>0$, in the side $N\ge4$ and in
$\mu\in\mathbb{R}$.
\begin{enumerate}
\item\label{it:all} For every $N\ge4$, if $\beta t^{2}/U\le1/6000$ then every pair
$x\ne y$ in $\Lambda$ satisfies
\begin{equation}\label{eq:main}
\bigl|\langle\Dop_x^\dagger\Dop_y\rangle\bigr|
\le\frac{57}{500}\left(\frac{99}{100}\right)^{d}
\qquad\text{with } d=d(x,y)\ge1,
\end{equation}
and consequently the rows of the correlation matrix are bounded by
$1+4A q/(1-q)^{2}\le4515.4$ uniformly in $N$, in $\mu$ and in $\Omega$, with
$A=57/500$ and $q=99/100$. Thus $k_{\mathrm{B}}T_{\mathrm{pair}}\le6000\,t^{2}/U$.
\item\label{it:even} If $N$ is even, the same conclusion \eqref{eq:main} holds under
the weaker hypothesis $\beta t^{2}/U\le1/4500$, hence
$k_{\mathrm{B}}T_{\mathrm{pair}}\le4500\,t^{2}/U$ on even tori.
\end{enumerate}
\end{theorem}

The two parts of the theorem carry the same prefactor and the same decay ratio, so
the difference between them is only the admissible temperature window. The
hypothesis is a condition on $\beta t^{2}/U$ alone: it is satisfied at high
temperature, it places no bound on $t$, $U$ or their ratio at fixed
$\beta t^{2}/U$, and it does not involve the phonon frequency, which may be sent to
zero or to infinity at fixed $U$ and $\Omega$ only through $K$ and $m$ with
$F^{2}/K$ and $\hbar^{2}K/m$ fixed. In the adiabatic limit $\Omega\to0$ the bound is
unchanged.

The proof occupies Sections~\ref{sec:bond} to~\ref{sec:decay}; it is assembled in
Section~\ref{sec:proof}. The ingredients are, in the order in which they are used:
the exact signless majorant of the oscillator contraction
(Section~\ref{sec:langfirsov}), the connected-support expansion of the correlation
(Section~\ref{sec:supports}), the exact sixteen-dimensional single-bond spectrum
(Section~\ref{sec:bond}), the bridge parity with its colouring and H\"older
decoupling (Section~\ref{sec:colour}), the counting of connected supports
(Section~\ref{sec:counting}), and the Koteck\'y--Preiss vertex criterion with
$B=557/500$ (Section~\ref{sec:kp}).

\subsection{What the constants are and where they come from}\label{sec:constants}

The three constants of Theorem~\ref{thm:main} are
\begin{equation}\label{eq:constants}
C=6000 \;(4500 \text{ on even tori}),\qquad q=\frac{99}{100},\qquad
A=\frac{57}{500}=0.114,
\end{equation}
with $k_{\mathrm{B}}T_{\mathrm{pair}}\le C\,t^{2}/U$. Internally the proof uses one
more pair of numbers, the radius $B=557/500=1.114$ and the insertion temperature
$T_0=303/250=1.212$ (and $121/100$ on even tori), whose powers generate the
coefficients of the cluster expansion; $A=B-1$ is the prefactor that the
Koteck\'y--Preiss theorem delivers, and $\ln B\ge0.1079571$ is the margin that the
budgets must fit under. Table~\ref{tab:audit} records the certified budget.

\begin{table}[t]
\centering
\small
\caption{The two certified budgets of Theorem~\ref{thm:main}. The row
$S_{\mathrm{phys}}$ splits into the twenty class slots of Table~\ref{tab:classes},
which contribute $0.0053568$ on all tori and $0.0071407$ on even tori, and the
$k\ge7$ tail of \eqref{eq:raney}, which contributes $0.0990539$ and $0.0948558$. All
entries are outward-rounded decimals of exact rational computations; the fourth row
is the ceiling $\ln B$ that they must stay under.}
\label{tab:audit}
\begin{tabular}{lcc}
\hline
quantity & all $N\ge4$, $C=6000$ & even $N$, $C=4500$\\
\hline
$S_{\mathrm{phys}}$ (physical supports) & $\le0.1044108$ & $\le0.1019965$\\
$S_{\mathrm{src}}$ (source supports) & $\le0.0008575$ & $\le0.0008381$\\
$S_{\mathrm{phys}}+S_{\mathrm{src}}$ & $\le0.1052682$ & $\le0.1028345$\\
$\ln B$ with $B=557/500$ & $\ge0.1079571$ & $\ge0.1079571$\\
tail test $B M_t^{1/c_{\mathrm{tail}}}/\sqrt{Cy_t}$ vs.\ $w_{\mathrm{eff}}q$
& $0.1178692\le0.1178862$ & $0.1172984\le0.1173541$\\
\hline
\end{tabular}
\end{table}

The tail test in the last row is the binding constraint of the whole scheme: the
margin is $1.7\times10^{-5}$ in absolute terms. The certified failure thresholds sit
just below the working constants, at $C=5800$ on all tori and $C=4300$ on even tori,
where the same tail test fails. Within this scheme $6000$ and $4500$ are therefore
close to the smallest constants that the enumerated classes and the four-ary tail
allow; Section~\ref{sec:floors} discusses how much of that is a property of the
method rather than of the model.

\section{The earlier route in the same setting}\label{sec:earlier}

The same model was treated in a previous release \cite{rens2026a}, where exponential
decay was proved at $\beta t^{2}/U\le2^{-27}$ with decay ratio $q=5/64$ and
prefactor $240/59\approx4.07$. That release is unconditional in $t/U$ and is
therefore the direct point of comparison. The present paper keeps the overall
architecture of that argument, a signless trajectory majorant, a polymer expansion
over connected hopping supports, and a Koteck\'y--Preiss test, and replaces each
quantitative step. Table~\ref{tab:audit2} records the replacement and the factor
gained in $C$.

\begin{table}[t]
\centering
\small
\caption{Where the improvement in $C$ against the earlier release \cite{rens2026a} comes
from. The entries of the last column are of different kinds: two are exact ratios of
the constants quoted, one is the ratio of a bound to a value, and three are
attributions rather than separable factors. They do not multiply out to the net
factor of $22\,369$; see the text.}
\label{tab:audit2}
\begin{tabular}{L{4.1cm}L{3.5cm}L{3.2cm}L{2.4cm}}
\hline
step & earlier release \cite{rens2026a} & this paper & factor in $C$\\
\hline
hopping/trace bound & on-site sum of squares $32kg^{2}/U$ &
exact sixteen-state bond spectrum with the half-filling supremum of
Theorem~\ref{thm:bond} & about $25$ (attribution)\\
atomic fugacity ratio & $z_b/z_a\le5/4$ &
eliminated: $R_e(b;g;z)\le R_e(b;g;1)$ proved in
Theorem~\ref{thm:bond} & $(5/4)^{2}=1.5625$ (exact)\\
bridge hopping multiplicity & $m_e\ge1$ &
$m_e\ge2$ on bridges, $m=k+b$ (Lemma~\ref{lem:bridge}) & not separable\\
rooted counting of supports & Euler tours, $16^{k}$ &
exact torus enumeration at $k\le6$ plus four-ary Raney tail
(Lemma~\ref{lem:counting}) & about $5.6$ (attribution)\\
Koteck\'y--Preiss radius & $B=e<3$ & $B=557/500$ & $(3/1.114)^{2}\approx7.25$\\
insertion scale & fixed length $1024$ & class dependent
$\ell_E=\sqrt{Cy_E}$ & coupled to $R_e$\\
pair prefactor & $240/59\approx4.07$ &
$A=B-1=57/500=0.114$ & $35.68$ (exact)\\
\hline
\end{tabular}
\end{table}

The factors in the table are not separable and do not multiply to the net
improvement. The atomic fugacity ratio $(5/4)^{2}=1.5625$ enters $C$ directly, the
radius ratio $7.25$ compares a bound with a value, and $35.68$ is the exact
improvement in the prefactor $A$, not a factor in $C$. The two entries that do enter
$C$ give about $11.3$, so a residual factor of about $2000$ is shared between the
bond spectrum, the bridge parity and the enumeration; we do not attempt to apportion
it further.

Two of the replacements are compensating rather than one-sided, and this is where
the honest comparison lives. First, the decay ratio degrades from $5/64\approx0.078$
to $99/100$, that is the correlation length grows from about one lattice spacing to
about a hundred; the improvement in $C$ is bought partly by accepting a bound that is
informative only at distances of a few hundred sites once the prefactor is included.
Second, the earlier release reached $2^{27}$ with a two-parameter family of ancestors
that included the cubic lattice; the present argument is genuinely
two-dimensional, through the four- or five-edge colouring and the four-ary tree tail,
and its transfer to higher dimensions is not part of this paper.

The companion preprint \cite{rens2026b} proves a dimension-tracking statement for the
same model on hypercubic tori, with
$C(d)=442368\,d^{2}(d+1)^{2}$ and hence $C(2)=15\,925\,248$, decay ratio
$q(2)=5/24$ and prefactor $A(2)=360/19$. In two dimensions the ceiling of this paper
is smaller by a factor of $2654$ on all tori and $3539$ on even tori at the same
ensemble, while its decay ratio is larger, $99/100$ against $5/24$. Neither constant
dominates the other: the dimension-tracking constant loses by a factor
$d^{2}(d+1)^{2}$ in dimension $d$, and the two-dimensional constant loses by a factor
$4.75$ in the ratio $q$, which at $q\approx1$ is a factor of $156$ in the exponential
rate $|\ln q|$ and hence in the correlation length.

\section{The oscillator contraction is a contraction}\label{sec:langfirsov}

The oscillator coordinate couples to the electronic density only through
$n_x-1$, so the on-site Lang--Firsov unitary \cite{hohenadler2007}
\[
\mathcal{U}=\exp\!\left(-\alpha_{\mathrm{LF}}\sum_{x\in\Lambda}(n_x-1)
(b_x^{\dagger}-b_x)\right),\qquad
\alpha_{\mathrm{LF}}=\frac{F}{\sqrt{2m\hbar\Omega^{3}}}
=\sqrt{\frac{U}{2\hbar\Omega}},
\]
with $X_x=\sqrt{\hbar/(2m\Omega)}\,(b_x+b_x^{\dagger})$, removes the linear
electron-phonon coupling and dresses the hopping with displacement operators
$D_x=\exp(-\alpha_{\mathrm{LF}}(b_x^{\dagger}-b_x))$ of unit operator norm. Because
the pair field $\Dop_x$ carries charge two, the displacement survives as
$D_x^{2}$ on a pair insertion and still has operator norm one.

Expanding $e^{-\beta H}$ in a time-ordered Dyson series in the interaction picture of
the decoupled Hamiltonian and integrating out the oscillators exactly gives, for a
closed electronic trajectory $\omega$ on $[0,\beta)$ with occupation
$n_x(u)\in\{0,1,2\}$ on each site, the normalized phonon weight
\begin{equation}\label{eq:phi}
\exp\!\left(-\sum_{x\in\Lambda}\Phi_\beta[n_x]\right),\qquad
\Phi_\beta[n]=\frac{\beta U}{2}\sum_{s\ne0}
\frac{\omega_s^{2}}{\omega_s^{2}+\Omega^{2}}\,|n_s|^{2}\ \ge0,
\end{equation}
the sum running over the occupation modes $n_s=\beta^{-1}\int_0^{\beta}
n(u)e^{-i\omega_su}du$ of the site, $\omega_s=2\pi s/\beta$, in units $\hbar=1$.
Appendix~\ref{app:oscillator} records the equivalent kernel form, in which
$\Phi_\beta[n]$ is a quarter of the double integral of $(n(u)-n(v))^{2}$ against the
positive kernel $\tfrac{U\Omega}{2}(e^{-\Omega w}+e^{-\Omega(\beta-w)})
/(1-e^{-\beta\Omega})$, and derives \eqref{eq:phi} from it. The derivation is local
in the site index; the positivity $\Phi_\beta[n]\ge0$ and the resulting bound
\begin{equation}\label{eq:majorant}
0\le\exp\!\left(-\sum_{x\in\Lambda}\Phi_\beta[n_x]\right)\le1
\end{equation}
hold verbatim in any dimension and in particular on the square torus. The trace
identity behind them is proved in \cite[Section 3]{rens2026a} and is repeated in
Appendix~\ref{app:oscillator} only to fix the notation. Since no sign of an
electronic matrix element, of a Jordan--Wigner
string or of a hopping amplitude survives \eqref{eq:majorant}, every trace below may
be replaced by the corresponding signless trace, in which each legal directed hop
contributes a matrix element one instead of a signed fermionic amplitude. The
signless trace is a bookkeeping majorant, not a Hamiltonian.

\section{The signless majorant and the connected-support expansion}
\label{sec:supports}

Rescale the hopping as $t=t/\ell_e\cdot\ell_e$ on every edge and let the edge scales
$\ell_e\ge1$ be free for the moment. Summing \eqref{eq:majorant} over trajectories
gives, for the vacuum polymers of the signless expansion, activities $w(E)$ supported
on connected edge sets $E\subset E(\Lambda)$ with $|E|\ge1$, and for the pair-marked
polymers activities $w_{xy}(E)$ on the same supports with $x,y\in V(E)$. With vertex
compatibility $E_1\sim E_2$ if and only if $V(E_1)\cap V(E_2)\ne\emptyset$, the
correlation admits the connected-support representation
\begin{equation}\label{eq:rep}
\langle\Dop_x^\dagger\Dop_y\rangle_{\mathrm{GC}}
=\sum_{\substack{E\subset E(\Lambda)\ \text{connected}\\ x,y\in V(E)}}
w_{xy}(E)\,\frac{Z_{\Lambda\setminus V(E)}}{Z_\Lambda},
\qquad x\ne y,
\end{equation}
where $Z_\Lambda$ is the polymer partition function of the connected hopping supports
and $Z_{\Lambda\setminus V(E)}$ the same partition function with the vertices of
$E$ deleted. Equation~\eqref{eq:rep} is the representation proved in
\cite[Section 5]{rens2026a}; nothing in this paper modifies it, and the entire content of
the present argument is a sharper bound on $|w_{xy}(E)|$ and a sharper summation over
supports.

The two gains of this paper enter \eqref{eq:rep} at different places. The bridge
parity of Section~\ref{sec:colour} counts the number of hops across every edge of a
support, so that an edge with many hops pays $\ell_e^{-m_e}$ and an edge that is a
bridge pays at least its square. The single-bond spectrum of Section~\ref{sec:bond}
bounds the signless trace that remains after the hops have been counted, and it is
uniform in the atomic parameter. Neither step uses smallness of $t/U$.

\section{The single bond: exact spectrum, half-filling supremum, certified
intervals}\label{sec:bond}

All the analytic input about the atoms is contained in the trace of the signless
hopping on one bond. On $e=\{u,v\}$ with Fock space $\mathcal{F}(e)$ of dimension
$16$, let $b=\beta U\ge0$, $g\ge0$, $z=e^{\beta\mu_c}>0$,
$q=e^{-b/2}\in(0,1]$ and $r=\sqrt{b^{2}/4+4g^{2}}\ge b/2$, and set
\begin{equation}\label{eq:Re}
R_e(b;g;z)=\frac{\Tr_{\mathcal{F}(e)}\bigl[
\exp(-\beta D_{0,e}+gB_e)\bigr]}{Z_0(e)},\qquad
Z_0(e)=\Tr_{\mathcal{F}(e)}e^{-\beta D_{0,e}} .
\end{equation}

\begin{theorem}[Exact five-sector spectrum and the half-filling
supremum]\label{thm:bond}
In the notation above:
\begin{enumerate}
\item\label{it:spec} The eigenvalues of $-\tfrac{b}{2}\NoddE+gB_e$, sorted by
total electron number $N_{\mathrm{tot}}\in\{0,1,2,3,4\}$, are
\[
\begin{aligned}
N_{\mathrm{tot}}=0 &: \{0\}\ \text{(mult.\ 1)},\\
N_{\mathrm{tot}}=1 &: \{-\tfrac{b}{2}+g\ \text{(mult.\ 2)},\;
-\tfrac{b}{2}-g\ \text{(mult.\ 2)}\},\\
N_{\mathrm{tot}}=2 &: \{-b\ \text{(mult.\ 3)},\; 0\ \text{(mult.\ 1)},\;
-\tfrac{b}{2}+r\ \text{(mult.\ 1)},\; -\tfrac{b}{2}-r\ \text{(mult.\ 1)}\},\\
N_{\mathrm{tot}}=3 &: \{-\tfrac{b}{2}+g\ \text{(mult.\ 2)},\;
-\tfrac{b}{2}-g\ \text{(mult.\ 2)}\},\\
N_{\mathrm{tot}}=4 &: \{0\}\ \text{(mult.\ 1)}.
\end{aligned}
\]
\item\label{it:closed} With $w=z+z^{-1}\ge2$, $C_0=\cosh(g)$ and
$C_2=\cosh(r)\ge\cosh(b/2)=\tfrac12(q+q^{-1})$,
\begin{equation}\label{eq:closed}
R_e(b;g;z)=1+\frac{A(b,g)\,w+B_0(b,g)}{(w+2q)^{2}},
\end{equation}
where $A(b,g)=4q(C_0-1)\ge0$ and $B_0(b,g)=2qC_2-q^{2}-1\ge0$.
\item\label{it:sup} For every $z>0$ and $b,g\ge0$,
\begin{equation}\label{eq:sup}
\begin{aligned}
R_e(b;g;z)&\le R_e(b;g;1)
=\frac{3+3q^{2}+8q\cosh(g)+2q\cosh(r)}{4(1+q)^{2}}\\
&\le\frac{3+4q^{2}+8q\cosh(g)+e^{r-b/2}}{4(1+q)^{2}}
=:R_e^{\mathrm{ub}}(b;g).
\end{aligned}
\end{equation}
\end{enumerate}
\end{theorem}

\begin{proof}
The Jordan--Wigner sign gauge of \cite[Section 4]{rens2026a} conjugates the physical
hopping on $e$ to $B_e$, which preserves $N_{\mathrm{tot}}$ and both spin numbers and
is the two-site signless hopping. On $N_{\mathrm{tot}}=0$ and $N_{\mathrm{tot}}=4$ the
space is one dimensional, $B_e=0$ and $\NoddE=0$. On $N_{\mathrm{tot}}=1$ and
$N_{\mathrm{tot}}=3$ the space is four dimensional, $\NoddE=1$, and $B_e$ is two
copies of $\bigl(\begin{smallmatrix}0&1\\1&0\end{smallmatrix}\bigr)$, giving
$-\tfrac{b}{2}\pm g$ twice each. On $N_{\mathrm{tot}}=2$ the two equal-spin states
give eigenvalue $-b$ with multiplicity two; on the remaining four-dimensional space
spanned by $|2,0\rangle,|0,2\rangle,|\uparrow\downarrow\rangle,|\downarrow\uparrow\rangle$
the matrix of $-\tfrac{b}{2}\NoddE+gB_e$ is
\[
\begin{pmatrix}0&0&g&g\\0&0&g&g\\g&g&-b&0\\g&g&0&-b\end{pmatrix},
\]
whose antisymmetric combinations give $0$ and $-b$ and whose symmetric combinations
span $\bigl(\begin{smallmatrix}0&2g\\2g&-b\end{smallmatrix}\bigr)$ with eigenvalues
$-\tfrac{b}{2}\pm r$. This proves the first item; the multiplicity count is
$1+4+6+4+1=16$.

For the second item, summing $z^{N_{\mathrm{tot}}}e^{\lambda}$ over the sixteen
eigenvalues gives
$Z_e=(1+z^{4})+(z+z^{3})\,4q\cosh(g)+z^{2}(1+3q^{2}+2q\cosh(r))$. Dividing by
$z^{2}$, using $z^{2}+z^{-2}=w^{2}-2$ and $Z_0(e)/z^{2}=(w+2q)^{2}$, gives the exact
polynomial identity
\[
\frac{Z_e(b;g;z)}{z^{2}}=(w+2q)^{2}+4q(C_0-1)w+(2qC_2-q^{2}-1),
\]
which is \eqref{eq:closed}; $A\ge0$ is immediate and $B_0\ge0$ follows from
$C_2\ge\cosh(b/2)$. For the third item, differentiate \eqref{eq:closed} in
$w\ge2$:
\[
\frac{\partial R_e}{\partial w}
=\frac{A(w+2q)-2(Aw+B_0)}{(w+2q)^{3}}
=-\frac{A(w-2q)+2B_0}{(w+2q)^{3}}\le0,
\]
since $w\ge2\ge2q$ and $A,B_0\ge0$. The supremum in $z>0$ is therefore attained at
$z=1$, that is $w=2$, which is the first inequality of \eqref{eq:sup}; the second
follows from $2q\cosh(r)=e^{r-b/2}+e^{-(r+b/2)}$ and
$e^{-(r+b/2)}\le e^{-b}=q^{2}$, so that
$3q^{2}+e^{-(r+b/2)}\le4q^{2}$.
\end{proof}

Item~\ref{it:sup} is what removes the atomic fugacity ratio $z_b/z_a\le5/4$ that the
earlier release had to carry: the supremum in the chemical potential is attained at
half filling and is available in closed form, so no ratio enters the budget.

What remains is a one-variable problem: bound
$\sup_{b\ge0}R_e^{\mathrm{ub}}(b;c\sqrt{yb})$ for the specific rational parameters
$(c,y)$ that the classes of Section~\ref{sec:counting} generate. Write
\begin{equation}\label{eq:Pe}
P_e(c,y):=\sup_{b\ge0}R_e\bigl(b;c\sqrt{yb};1\bigr)\le\sup_{b\ge0}
R_e^{\mathrm{ub}}\bigl(b;c\sqrt{yb}\bigr),
\end{equation}
where the argument $g=c\sqrt{yb}$ is the only place the colour count $c$ and the
scaled hopping enter. The bound $g=c\beta\ell_Et\le c\sqrt{y_E b}$ used here is the
statement that $\beta\ell_Et\le\sqrt{y_E b}$, which holds whenever $\ell_E=\sqrt{Cy_E}$ and
$\beta t^{2}/U\le1/C$, since then
$\beta\ell_Et=\sqrt{y_EC(\beta t^{2}/U)b}\le\sqrt{y_Eb}$.

\begin{corollary}[Certified uniform bounds]\label{cor:pe}
For every $C>0$ and every parameter pair $(c,y)$ appearing in
Tables~\ref{tab:classes} and~\eqref{eq:raney}, the supremum \eqref{eq:Pe} is
bounded above by an explicit rational $M=M(c,y)$, uniformly in $b\ge0$ and in
$\beta t^{2}/U\le1/C$. The twenty-two class slots of the two tables use twelve
distinct rational triples $(c,y,M)$, and all of them are certified in exact
outward-rounded rational arithmetic on $102$ adaptive
intervals covering $[0,40]$, together with an exact tail bound on $[40,\infty)$
obtained from $r-b/2=4c^{2}yb/(r+b/2)\le4c^{2}y$.
\end{corollary}

The tail estimate in the corollary is where the colour count becomes expensive:
$c$ enters the exponential through $c^{2}y$, and a larger $c$ forces a larger
$M$. This is the amplification discussed in Section~\ref{sec:floors}.

\section{Bridge parity, edge colouring, and coloured H\"older decoupling}
\label{sec:colour}

A support $E$ with $k$ edges and $b(E)$ bridges constrains how often a closed
trajectory can cross each edge. The first statement is a parity count.

\begin{lemma}[Bridge parity]\label{lem:bridge}
Let $E$ be a connected subgraph with $k=|E|\ge1$ edges and $b=b(E)$ bridges. For any
closed electronic trajectory $\omega$ with active edge set $E$, either the vacuum
trajectory contributing to $w(E)$ or a trajectory with one insertion
$O_{xy}=\Dop_x^\dagger\Dop_y$, $x\ne y\in V(E)$, contributing to $w_{xy}(E)$, let
$m_e(\omega)\ge1$ be the number of hops across $e\in E$. Then $m_e(\omega)$ is even,
hence at least two, on every bridge $e$, and at least one on every other edge; in
particular $\sum_{e\in E}m_e(\omega)\ge m(E):=k+b(E)$.
\end{lemma}

\begin{proof}
Let $e=\{u,v\}$ be a bridge of a physical component and let $V_1\sqcup V_2$ be the
vertex partition it defines, with $u\in V_1$. No edge of the component other than $e$
joins $V_1$ to its complement, so the electron number $N(V_1)$ changes only through
hops across $e$ and through pair insertions. Along the closed trajectory
$N(V_1)$ returns to its initial integer value. Each hop across $e$ changes $N(V_1)$ by
$\pm1$, that is by an odd amount. The insertion $\Dop_x^\dagger\Dop_y$ creates two
electrons at $x$ and annihilates two electrons at $y$, and each pair-source jump does
the same, so the change of $N(V_1)$ is $0$ or $\pm2$, always even. Reducing
$\Delta N(V_1)=0$ modulo $2$ gives $m_e\equiv0\pmod2$; since $e$ is active,
$m_e\ge2$. Summing over edges gives the stated total.
\end{proof}

The lemma has a consequence that is easy to miss: the interior of a large
one-dimensional chain pays the same squared factor on each of its links, so a chain
support of $k$ edges pays $\ell^{-2k}$ rather than $\ell^{-k}$.

\begin{lemma}[Edge colouring]\label{lem:edgecolor}
Let $T_N=(\mathbb{Z}/N\mathbb{Z})^{2}$.
\begin{enumerate}
\item\label{it:chi} If $N\ge4$ is even, $\chi'(T_N)=4$, explicitly: colour the
horizontal edge $\{(x,y),(x+1,y)\}$ by $1$ or $2$ according to $x\bmod2$, and the
vertical edge $\{(x,y),(x,y+1)\}$ by $3$ or $4$ according to $y\bmod2$.
\item\label{it:chi5} If $N\ge5$ is odd, $\chi'(T_N)=5$, explicitly: colour the
horizontal edges $\{(x,y),(x+1,y)\}$ with even $x\in\{0,2,\dots,N-3\}$ together with
the vertical wrap-around edge $\{(N-1,N-1),(N-1,0)\}$ by $1$; the horizontal edges
with odd $x\in\{1,3,\dots,N-2\}$ together with $\{(0,N-1),(0,0)\}$ by $2$; the
vertical edges $\{(x,y),(x,y+1)\}$ with even $y\in\{0,2,\dots,N-3\}$ by $3$; those
with odd $y\in\{1,3,\dots,N-2\}$ by $4$; and all remaining wrap-around edges (the
horizontal edges $\{(N-1,y),(0,y)\}$ and the vertical edges $\{(x,N-1),(x,0)\}$ for
$1\le x\le N-2$) by $5$.
\item\label{it:k6} For every connected subgraph $E\subset T_N$ with $N\ge4$ and
$1\le|E|\le6$, one has $\chi'(E)\le c_{\mathrm{ub}}(E)$, where
$c_{\mathrm{ub}}(E)=\Delta(E)$ if $E$ is bipartite and
$c_{\mathrm{ub}}(E)=\Delta(E)+1$ otherwise; in particular $c_{\mathrm{ub}}(E)\le4$
on every class of Table~\ref{tab:classes}.
\end{enumerate}
\end{lemma}

\begin{proof}
For (1), every vertex has one horizontal edge of each parity and one vertical edge of
each parity, so the four classes are matchings and the four-colouring is proper; a
$4$-regular graph needs at least four colours. For (2), let $(x,y)$ be a vertex. If
$1\le x\le N-2$ the two incident horizontal edges carry colours $1$ and $2$, the two
non wrap-around vertical edges carry $3$ and $4$, and a vertical wrap-around edge,
present only when $y\in\{0,N-1\}$, carries $5$. If $x=0$ the incident horizontal
edges are $\{(0,y),(1,y)\}$ of colour $1$ and $\{(N-1,y),(0,y)\}$ of colour $5$, the
vertical edges carry $3$ and $4$, and the vertical wrap-around edge
$\{(0,N-1),(0,0)\}$ carries $2$. If $x=N-1$ the incident horizontal edges carry $2$
and $5$, the vertical edges carry $3$ and $4$, and the vertical wrap-around edge
$\{(N-1,N-1),(N-1,0)\}$ carries $1$. In each case the incident colours are distinct,
so the five colouring is proper. It is optimal because a matching on the
$N^{2}$ vertices with $N$ odd has at most $(N^{2}-1)/2$ edges, so four colours cover
at most $2(N^{2}-1)<2N^{2}=|E(T_N)|$ edges; hence $\chi'(T_N)=5$. For (3), K\"onig's
line-colouring theorem gives $\chi'(E)=\Delta(E)$ on bipartite graphs and Vizing's
theorem gives $\chi'(E)\le\Delta(E)+1$ in general. A connected subgraph with at most
six edges on a square torus is non-bipartite only if it contains an odd cycle, and the
only odd cycles available are the two rooted five-cycles and the five-cycle with one
adjacent edge; a five-cycle has $\Delta=2$ and the latter $\Delta=3$, so
$c_{\mathrm{ub}}\le4$. Every class of Table~\ref{tab:classes} is checked in that
table.
\end{proof}

Part~\ref{it:k6} is the only place where the small-edge enumeration pays off in the
colouring: for $k\ge7$ the supports may contain longer odd cycles and up to four
edges at a vertex, and the colour count is then $c_{\mathrm{tail}}=5$ on all tori and
$4$ on even tori, where the grid is bipartite.

\begin{lemma}[Coloured H\"older decoupling]\label{lem:holder}
Let $E$ be a finite subgraph with a proper edge colouring
$E=\bigsqcup_{j=1}^{c}E_j$ into matchings. For edge scales $\ell_e\ge1$ let
$H_j=-D_{0,V(E)}+c\sum_{e\in E_j}\ell_etB_e$ and
$Q_E=\exp\bigl(\beta(-D_{0,V(E)}+\sum_{e\in E}\ell_etB_e)\bigr)$. Then
\begin{equation}\label{eq:holder}
\begin{aligned}
\Tr_{\mathcal{F}(V(E))}\Bigl[\exp\Bigl(-\beta D_{0,V(E)}
+\beta\sum_{e\in E}\ell_etB_e\Bigr)\Bigr]
&\le\prod_{j=1}^{c}\Bigl(\Tr_{\mathcal{F}(V(E))}e^{\beta H_j}\Bigr)^{1/c}\\
&=Z_0(V(E))\prod_{e\in E}R_e\bigl(b;\,c\ell_e\beta t;\,z\bigr)^{1/c},
\end{aligned}
\end{equation}
and for the pair-marked trace with $O_{xy}=\Dop_x^\dagger\Dop_y$, $x\ne y$,
\[
\bigl|\Tr_{\mathcal{F}(V(E))}\bigl(O_{xy}Q_E\bigr)\bigr|
\le\Tr_{\mathcal{F}(V(E))}\bigl(Q_E\bigr),
\]
so that \eqref{eq:holder} applies to $|w_{xy}(E)|$ as well.
\end{lemma}

\begin{proof}
The pair-marked statement is immediate: $Q_E$ is Hermitian positive semidefinite with
eigenvalues $\lambda_a\ge0$ and $O_{xy}$ is a partial isometry of operator norm one,
so $|\Tr(O_{xy}Q_E)|\le\sum_a\lambda_a=|\Tr Q_E|$. For the first statement, the
Lie--Trotter product formula \cite{trotter1959} applied to
$\beta\sum_jH_j/c$ and the Schatten H\"older inequality for trace ideals give, for
every finite product length $M$, the bound
$|\Tr\prod_{m}\prod_{j}A_j|\le\prod_{j}(\Tr A_j^{cM})^{1/c}$, and letting $M\to\infty$
yields the first inequality of \eqref{eq:holder}. Within each colour class the edges
are vertex disjoint, so $H_j$ is a direct sum of single-bond terms and the
trace factorizes: $\Tr e^{\beta H_j}/Z_0(V(E))=\prod_{e\in E_j}R_e(b;c\ell_e\beta
t;z)$. Multiplying over the colour classes gives the identity in
\eqref{eq:holder}.
\end{proof}

The price of the decoupling is visible in the argument: the bond ratio is evaluated
at $c\ell_e\beta t$ rather than $\ell_e\beta t$, so the hopping appears amplified by
the colour count. That factor is the origin of the $c^{2}$ in
Section~\ref{sec:bond} and, quantitatively, of the floor in
Section~\ref{sec:floors}.

\section{Counting connected supports}\label{sec:counting}

Two counting devices are used. Supports with at most six edges are enumerated exactly
on the torus, because the single-bond bound is not enough there and the bridge parity
is most valuable; larger supports are counted by a tree injection.

\begin{lemma}[Exact classes and the Raney tail]\label{lem:counting}
Let $T_N$ be a square torus with $N\ge4$.
\begin{enumerate}
\item\label{it:exact} Every connected edge subgraph of $T_N$ with $1\le|E|\le6$
belongs, together with its root vertex, to one of the twenty classes of
Table~\ref{tab:classes}, which record the number of vertices $v$, the number of edges
$k$, the colour count $c=c_{\mathrm{ub}}$, the hop total $m=k+b(E)$ and the
multiplicity. On an even torus the two classes marked with a star do not occur, so
eighteen classes suffice.
\item\label{it:tail} On a four-regular graph, an ordered depth-first exploration from
a fixed root with four root slots and three forward slots per traversed edge maps
every connected $k$-edge subgraph that contains the root injectively to a rooted
quadruple of trees of total edge count $k$. The generating function of such
quadruples is $T(w)^{4}$ with $T=1+wT^{3}$, whose coefficients are
\begin{equation}\label{eq:raney}
R_4(k)=[w^{k}]T(w)^{4}=\frac{4}{3k+4}\binom{3k+4}{k},\qquad
R_4(0,\dots,6)=(1,4,18,88,455,2448,13566).
\end{equation}
Consequently, for any $w_{\mathrm{eff}}\in(0,4/27)$, with $T\in(1,3/2)$ the positive
root of $T=1+w_{\mathrm{eff}}T^{3}$,
\[
\sum_{\substack{E\ \text{connected},\,x\in V(E)\\ |E|\ge7}}
w_{\mathrm{eff}}^{|E|}\le T^{4}-\sum_{k=0}^{6}R_4(k)w_{\mathrm{eff}}^{k}.
\]
\item\label{it:twovar} If each site also carries two source incidences, the ordered
exploration on the resulting $(4,2)$-regular multigraph has generating function
\begin{equation}\label{eq:twovar}
G_{\mathrm{aug}}(w_p,w_s)=P^{4}S^{2},\qquad
P=1+w_pP^{3}S^{2},\qquad S=1+w_sP^{4}S=\frac{1}{1-w_sP^{4}},
\end{equation}
and every connected polymer $E=E_{\mathrm{phys}}\sqcup E_{\mathrm{src}}$ containing
the root with $k_p\ge0$ physical and $k_s\ge1$ source edges is counted at most
$[w_p^{k_p}w_s^{k_s}]G_{\mathrm{aug}}$. Moreover, for $T_0>1$,
$w_{\mathrm{eff}}=(T_0-1)/T_0^{3}$, $w_sT_0^{4}<1$ and
$w_p\le w_{\mathrm{eff}}/S_{\mathrm{hi}}^{2}$ with
$S_{\mathrm{hi}}=(1-w_sT_0^{4})^{-1}$, the monotone fixed point satisfies
$P\le T_0$ and $S\le S_{\mathrm{hi}}$, and
\begin{equation}\label{eq:twobound}
G_{\mathrm{aug}}(w_p,w_s)-G_{\mathrm{aug}}(w_p,0)
\le w_s\bigl(4T_0^{3}S_{\mathrm{hi}}^{2}D_P+2T_0^{4}S_{\mathrm{hi}}D_S\bigr),
\end{equation}
with $D_P=2w_pT_0^{7}S_{\mathrm{hi}}^{3}/\mathrm{den}$,
$D_S=T_0^{4}S_{\mathrm{hi}}^{2}+4w_sT_0^{3}S_{\mathrm{hi}}^{2}D_P$ and
$\mathrm{den}=1-3w_pT_0^{2}S_{\mathrm{hi}}^{2}-8w_pw_sT_0^{6}S_{\mathrm{hi}}^{3}>0$.
\end{enumerate}
\end{lemma}

\begin{proof}
For (1), exhaustive breadth-first enumeration of rooted connected edge subgraphs of
$T_N$ for $N=4,5,6,7$ produces the twenty classes; a connected subgraph with at most
six edges has diameter at most six, so every cycle in it has length at most six and
each such subgraph on $T_N$ with $N\ge7$ lifts isomorphically to the infinite grid,
which makes $T_7$ exhaustive for all $N\ge7$. On an even torus the grid is bipartite,
so neither the five-cycle nor the five-cycle with an edge appears. For (2), fix an
ordering of the four physical incidences at every vertex and run depth-first search.
The root opens four slots; traversing an edge into an unvisited vertex leaves three
forward slots, which is the relation $T=1+wT^{3}$; an edge whose far endpoint is
already visited closes a cycle and is recorded as a terminal leaf, which is allowed
because $T\ge1$ termwise. Since every edge is discovered from its earlier endpoint
and the incidence ordering identifies the edge, the map to marked trees is injective.
Lagrange inversion gives the coefficient formula in \eqref{eq:raney}. For (3), the
same argument with two source incidences per site gives \eqref{eq:twovar}; bounding
$G_{\mathrm{aug}}(w_p,w_s)-G_{\mathrm{aug}}(w_p,0)$ by the mean value theorem in
$w_s$ and differentiating the implicit equations under $P\le T_0$, $S\le S_{\mathrm{hi}}$
gives \eqref{eq:twobound}.
\end{proof}

\begin{table}[t]
\centering
\footnotesize
\caption{The twenty classes of connected supports with at most six edges, by the
number of vertices $v$, the number of edges $k$, the colour count $c$, the hop total
$m=k+b$ and the multiplicity. A star marks the two classes that cannot occur on an
even torus. Every entry is enumerated exactly; the last column is the contribution of
the class to the budget $S_{\mathrm{phys}}$ at $C=6000$, computed in exact rational
arithmetic and rounded to nearest at seven decimals, so the exact sum over the twenty
slots, $0.0053568$, differs from the sum of the printed entries by the rounding.}
\label{tab:classes}
\begin{tabular}{rrrrr r}
\hline
$v$ & $k$ & $c$ & $m$ & multiplicity & contribution at $C=6000$\\
\hline
2 & 1 & 1 & 2 & 4 & $0.0047994$\\
3 & 2 & 2 & 4 & 18 & $0.0001099$\\
4 & 3 & 2 & 6 & 72 & $0.0000011$\\
4 & 3 & 3 & 6 & 16 & $0.0000013$\\
4 & 4 & 2 & 4 & 6 & $0.0001061$\\
5 & 4 & 2 & 8 & 250 & $0.0000045$\\
5 & 4 & 3 & 8 & 180 & $0.0000032$\\
5 & 4 & 4 & 8 & 5 & $0.0000001$\\
5 & 5 & 3 & 5 & $2^{\star}$ & $0.0000087$\\
5 & 5 & 3 & 6 & 60 & $0.0000141$\\
6 & 5 & 2 & 10 & 852 & $0.0000171$\\
6 & 5 & 3 & 10 & 1260 & $0.0000253$\\
6 & 5 & 4 & 10 & 72 & $0.0000014$\\
6 & 6 & 2 & 6 & 48 & $0.0000039$\\
6 & 6 & 3 & 8 & 432 & $0.0000088$\\
6 & 6 & 4 & 7 & $24^{\star}$ & $0.0000019$\\
6 & 6 & 4 & 8 & 36 & $0.0000007$\\
7 & 6 & 2 & 12 & 2730 & $0.0000617$\\
7 & 6 & 3 & 12 & 7658 & $0.0001732$\\
7 & 6 & 4 & 12 & 630 & $0.0000142$\\
\hline
\end{tabular}
\end{table}

The classes are cheap; the single bond, counted four times from a root, contributes
$0.0047994$, a twentieth of the whole physical budget, and the largest of the
remaining contributions is the six-edge, twelve-hop class $(7,6,3,12)$ at
$0.0001732$. The budget is dominated by the tail, which is where the four-ary tree
count with growth rate $1/w_{\mathrm{eff}}=8.4$ per edge replaces the Euler-tour
count with growth rate $16$ per edge of the earlier release.

\section{The Koteck\'y--Preiss vertex criterion}\label{sec:kp}

We quote the criterion in the form in which it is used.

\begin{proposition}[Koteck\'y--Preiss \cite{kotecky1986}, Theorem 1]\label{prop:kp}
Let $\mathcal{P}$ be a countable polymer set with the symmetric incompatibility
relation $\gamma\sim\gamma'$ if and only if $V(\gamma)\cap V(\gamma')\ne\emptyset$ and
complex activities $w(\gamma)$. Suppose that there are
$a,b\colon\mathcal{P}\to[0,\infty)$ such that for every $\gamma_0\in\mathcal{P}$
\begin{equation}\label{eq:kp}
\sum_{\gamma\sim\gamma_0}|w(\gamma)|\,e^{a(\gamma)+b(\gamma)}\le a(\gamma_0).
\end{equation}
Then the cluster expansion of $\ln Z_\Lambda$ converges absolutely, and for every
$\gamma_0$ one has
\[
\sum_{X\sim\gamma_0}|\Phi^{T}(X)|e^{b(X)}\le a(\gamma_0),
\]
with $b(X)=\sum_{\gamma\in X}b(\gamma)$. With $a(\gamma)=|V(\gamma)|\ln B$ and
$b(\gamma)=|E_{\mathrm{phys}}(\gamma)|\ln(1/q)$ for $B>1$ and $q\in(0,1)$, where
$E(\gamma)=E_{\mathrm{phys}}(\gamma)$ is the physical edge set of $\gamma$, the
vertex-sum condition
\begin{equation}\label{eq:vertexsum}
\sup_{x\in\Lambda}\sum_{\gamma\ni x}|w(\gamma)|\,B^{|V(\gamma)|}q^{-|E(\gamma)|}
\le\ln B
\end{equation}
implies \eqref{eq:kp}, and then for every $S\subset\Lambda$
\begin{equation}\label{eq:deletion}
\left|\frac{Z_{\Lambda\setminus S}}{Z_\Lambda}\right|
\le\exp\Bigl(\sum_{x\in S}\sum_{X\ni x}|\Phi^{T}(X)|\Bigr)\le B^{|S|}.
\end{equation}
\end{proposition}

The condition \eqref{eq:vertexsum} is used with $q=99/100$ and $B=557/500$, so the
right side is $\ln B\in[0.1079571,0.114]$, and with activities bounded in the form
\begin{equation}\label{eq:activity}
|w(E)|\le\ell_E^{-m(E)}P_e\bigl(c,y_E\bigr)^{k/c}
\le(Cy_E)^{-m(E)/2}M_E^{k/c}=:R_E,
\end{equation}
which combines Lemma~\ref{lem:bridge}, Lemma~\ref{lem:holder} and
Corollary~\ref{cor:pe}, and holds equally for $|w_{xy}(E)|$. Because
$A=B-1$, the prefactor delivered by \eqref{eq:deletion} is exactly the constant of
Theorem~\ref{thm:main}.

\section{Proof of the main theorem}\label{sec:proof}

The proof is a single chain of inequalities. The last two sections supply it; what
remains is bookkeeping.

\subsection{The physical vertex sum}

Fix $C$ and the pair $(B,T_0)$ as in Table~\ref{tab:audit}, and set
$w_{\mathrm{eff}}=(T_0-1)/T_0^{3}$. For a connected support $E$ with $k$ edges,
$v$ vertices, $b(E)$ bridges and colour count $c$, use the class-dependent scale
$\ell_E=\sqrt{Cy_E}\ge1$; the accompanying table certifies $Cy_E\ge1$ for every
class, which is what makes the scale admissible. By \eqref{eq:activity}, the
contribution of $E$ to the left side of \eqref{eq:vertexsum} is at most
$B^{v}R_Eq^{-k}$, and summing this over the classes of Table~\ref{tab:classes} and
over the tail gives
\begin{equation}\label{eq:physsum}
S_{\mathrm{phys}}=\sum_{k\le6\ \text{classes}}B^{v}R_Eq^{-k}
+\sum_{k\ge7}B^{|V(E)|}R_Eq^{-k}\le0.1044108
\end{equation}
on all tori at $C=6000$ and at most $0.1019965$ on even tori at $C=4500$. The tail is
not summed term by term: it is dominated by \eqref{eq:raney} at rate
$w_{\mathrm{eff}}$, and the resulting geometric sum is controlled by the single
inequality
\begin{equation}\label{eq:tailtest}
\frac{B\,M_t^{1/c_{\mathrm{tail}}}}{\sqrt{Cy_t}}
\le w_{\mathrm{eff}}\,q,
\end{equation}
with $(c_t,y_t,M_t)=(5,29/1000,662/125)$ on all tori and
$(4,39/1000,1893/500)$ on even tori. The left and right sides of \eqref{eq:tailtest}
are $0.1178692$ and $0.1178862$ on all tori and $0.1172984$ and $0.1173541$ on even
tori; these are the numbers in the last row of Table~\ref{tab:audit}, and the small
margin is the only place where the constants are tight.

\subsection{From the vertex sum to the correlation}

Every connected support in \eqref{eq:rep} that contains both $x$ and $y$ has at least
$d=d(x,y)$ edges, because a connected graph on $k$ edges has diameter at most $k$.
Hence $1\le q^{d-|E|}$ and, by \eqref{eq:rep}, \eqref{eq:deletion} and
\eqref{eq:vertexsum},
\begin{equation}\label{eq:chain}
\begin{aligned}
\bigl|\langle\Dop_x^\dagger\Dop_y\rangle\bigr|
&\le\sum_{\substack{E\ \text{conn.}\\ x,y\in V(E)}}|w_{xy}(E)|\,B^{|V(E)|}\\
&\le q^{d}\sum_{\substack{E\ \text{conn.}\\ x\in V(E)}}R_E\,B^{|V(E)|}q^{-|E|}
\le(\ln B)\,q^{d}\le(B-1)q^{d}=Aq^{d},
\end{aligned}
\end{equation}
the last step because $\ln B\le B-1$ and $A=B-1=57/500$. This proves
\eqref{eq:main}. The rows are then summable: on the torus there are at most $4d$ sites
at distance $d$, so the row sum at $x$ is at most
$1+\sum_{d\ge1}4d\,Aq^{d}=1+4Aq/(1-q)^{2}=4515.4$, uniformly in $N$; this is the
row-sum statement of Theorem~\ref{thm:main}\ref{it:all}. Since the argument
never uses the sign of $t$, of $\mu_c$ or of $F$, and the harmonic oscillators were
integrated out exactly at fixed $\Omega$, all constants are uniform in the
phonon frequency.

\subsection{Decay rate, prefactor, and the interpretation}\label{sec:decay}

Three features of the bound are worth separating. The exponential rate
$|\ln q|=0.01005$ is weak: the correlation length is about one hundred lattice
spacings, so the bound falls below one per cent only at distances of about two
hundred and forty sites, and it is useful as a statement about uniform summability
rather than as a quantitative decay estimate at moderate distance. The prefactor
$A=B-1=0.114$ is small, which is a consequence of the Koteck\'y--Preiss radius being
close to one; being below the trivial bound $1/2$ on the pair density, it makes the
bound informative at every distance. The constant $C$ in the temperature hypothesis
is the quantity that this paper improves by four orders of magnitude, and the product
$C\,t^{2}/U$ is the certified upper bound on $k_{\mathrm{B}}T_{\mathrm{pair}}$.

\section{The canonical thermodynamic limit}\label{sec:canonical}

The grand-canonical bound \eqref{eq:chain} does not by itself describe a system at
fixed density, because changing $N$ at fixed $\mu$ changes the density. The transfer
to the canonical ensemble is done, as in \cite[Section 6]{rens2026a}, by adding a
pair-transfer source
\begin{equation}\label{eq:source}
X(\xi)=\sum_{x\in\Lambda}\bigl[e^{i\theta}\Dop_x^\dagger\Dop_{x+r}
+\text{h.c.}\bigr],
\end{equation}
with a fixed lattice vector $r\ne0$, differentiating the source-augmented pressure in
$\xi$ on a complex disk, and controlling the derivative by a Cauchy estimate. The
source introduces polymers with source edges in addition to the physical supports,
and the counting of those is part~3 of Lemma~\ref{lem:counting}.

The source activity is bounded as follows. Separating the $k_s$ source insertions at
ordered times, bounding the phonon weight by one between them, and applying Schatten
H\"older to the commuting positive propagators with
$\|B_{e_s}\|_{\mathrm{op}}=1$ gives, for a polymer
$E=E_{\mathrm{phys}}\sqcup E_{\mathrm{src}}$ with physical components $E_1,\dots,E_r$
and $|V(E)|\le k_p+k_s+1$,
\begin{equation}\label{eq:srcbound}
\bigl|w(E_{\mathrm{phys}}\sqcup E_{\mathrm{src}})\bigr|
\le L_{\mathrm{src}}^{-k_s}
\bigl(e^{2\beta L_{\mathrm{src}}|\xi|}-1\bigr)^{k_s}\prod_{i=1}^{r}R_{E_i}.
\end{equation}
With $L_{\mathrm{src}}=10^{5}$ and $|\xi|\le\xi_0=\tau/(2L_{\mathrm{src}})$ one has
$e^{2\beta L_{\mathrm{src}}|\xi|}-1\le e-1<2<3$, so each source edge carries the
weight $w_s=3B/L_{\mathrm{src}}$. Bounding each physical component at
$B^{k_i}R_{E_i}q^{-k_i}\le(w_{\mathrm{eff}}/S_{\mathrm{hi}}^{2})^{k_i}$, which is the
component-wise form of \eqref{eq:physsum} and is checked class by class and for the
tail, the source part of the vertex sum is at most
\begin{equation}\label{eq:srcsum}
S_{\mathrm{src}}\le B\,w_s\,\partial_{w_s}G_{\mathrm{aug}}\le0.0008575
\quad(\text{all tori}),\qquad \le0.0008381\quad(\text{even tori}),
\end{equation}
the two values obtained from \eqref{eq:twobound}. Adding these to
\eqref{eq:physsum} gives the third row of Table~\ref{tab:audit},
$0.1052682\le0.1079571\le\ln B$ and $0.1028345\le0.1079571\le\ln B$, so the
vertex-sum condition holds on the whole disk $|\xi|\le\xi_0$, and the Cauchy estimate
transfers \eqref{eq:main} to translation-invariant thermodynamic canonical limits at
every interior density. This gives the following addition to
Theorem~\ref{thm:main}.

\begin{proposition}[Canonical transfer]\label{prop:canonical}
Under the hypotheses of Theorem~\ref{thm:main}\ref{it:all} (respectively
\ref{it:even}), with the same $(C,q,A)=(6000,99/100,57/500)$ (respectively
$(4500,99/100,57/500)$), bound \eqref{eq:main} holds for the canonical
thermodynamic limits at every interior total electron density $\rho\in(0,2)$,
including the balanced spin sectors $\rho_\uparrow=\rho_\downarrow=\rho/2$, uniformly
in $\Omega>0$.
\end{proposition}

The source scale $L_{\mathrm{src}}=10^{5}$ is a technical parameter: it is large
because the source derivative must be bounded on a small disk, and it costs only the
factor $w_s\approx0.0000334$ in \eqref{eq:srcsum}. The second and third rows of
Table~\ref{tab:audit} show the cost is genuinely small: the physical budget uses
$0.1044108$ of the available $0.1079571$ and the canonical transfer adds $0.0008575$.

\section{The trade-off curve}\label{sec:tradeoff}

The constants of Theorem~\ref{thm:main} are the favourable end of a one-parameter
family. The decay ratio $q$ and the prefactor $A=B-1$ both move with $B$, and the tail
test \eqref{eq:tailtest} is a comparison between $q$ and the growth rate of the
support count, so a faster decay, that is a smaller $q$, is affordable only at a
larger ceiling: at $q=1/2$ the certified ceiling rises to $26\,000$ on all tori and
$19\,800$ on even tori, and it decreases monotonically to $6000$ and $4500$ at
$q=99/100$. Theorem~\ref{thm:main} quotes the point of the family with the smallest
certified ceiling.

\begin{theorem}[Certified trade-off curve]\label{thm:curve}
For the same model and with no smallness hypothesis on $t/U$, the bound
$\bigl|\langle\Dop_x^\dagger\Dop_y\rangle\bigr|\le A_{\mathrm{pair}}q_{\mathrm{pair}}^{d}$
holds at the following certified rational points.
On all tori $N\ge4$:
\[
\begin{gathered}
(q_{\mathrm{pair}},C,A_{\mathrm{pair}})\in\bigl\{
(99/100,\,6000,\,57/500),\ (19/20,\,6550,\,53/500),\\
(9/10,\,7300,\,14/125),\ (3/4,\,10600,\,29/250),\
(1/2,\,26000,\,4/25)\bigr\},
\end{gathered}
\]
and on even tori $N\ge4$:
\[
\begin{gathered}
(q_{\mathrm{pair}},C,A_{\mathrm{pair}})\in\bigl\{
(99/100,\,4500,\,57/500),\ (19/20,\,4900,\,11/100),\\
(9/10,\,5450,\,57/500),\ (3/4,\,7900,\,59/500),\
(1/2,\,19800,\,41/250)\bigr\}.
\end{gathered}
\]
Only the two first points, with $C=6000$ and $C=4500$, also carry the canonical
transfer of Proposition~\ref{prop:canonical}; at the other eight the vertex sum is
certified without the source term.
\end{theorem}

The curve is the honest statement about what the scheme can do. Halving the decay
ratio from $99/100$ to $1/2$ costs a factor of four in the ceiling, from $6000$ to
$26000$ on all tori, and the prefactor rises from $0.114$ to $0.16$; the scheme
cannot buy a good decay ratio without a substantially weaker temperature bound.

\section{The weighted substrate route and its obstruction}\label{sec:weighted}

A different and older route to the same theorem splits the hopping off the atomic
Hamiltonian by a weighted two-parameter inequality, keeping a fraction of the atomic
gap and paying for the rest with a scalar term. The route has the advantage of being
close to elementary, and the disadvantage of carrying a smallness hypothesis on
$t/U$. It is useful to record exactly how much it can buy, because that is the
alternative to the exact bond spectrum used here.

\begin{theorem}[Weighted quadratic form and the $\varepsilon_0$
obstruction]\label{thm:weighted}
Let $E$ be a connected physical subgraph with $k=|E|$ edges and maximum degree four,
let $\gamma\in(0,1]$, $\kappa\in(0,1)$ and $g=\ell t\ge0$ satisfy
\begin{equation}\label{eq:kappa}
\frac{16g}{U}\le\kappa\gamma,
\qquad\text{i.e.}\qquad
\frac{t}{U}\le\varepsilon_0:=\frac{\kappa\gamma}{16\ell}.
\end{equation}
Then on $\mathcal{F}(V(E))$,
\begin{equation}\label{eq:twoab}
gB_E\le\frac{\gamma U}{4}\Nodd{}_{V(E)}
+\frac{32kg^{2}}{\gamma(1-\kappa)U}\,\mathbf{1}.
\end{equation}
Consequently, bounding the residual atomic fugacity factor at residual gap
$(1-\gamma/2)$ by $2$ and applying the bridge parity of Lemma~\ref{lem:bridge} with a
class-independent scale $\ell=\sqrt{C_{2ab}y}$, the vertex-sum condition holds at the
following five certified rational points, with $q=99/100$ and $B=113/100$ throughout
and with the source-free budget $S$ (last column, outward-rounded) compared to
$\ln B\ge0.122217$:
\begin{center}
\begin{tabular}{cccccc}
\hline
$\kappa$ & $\gamma$ & $y$ & $C_{2ab}$ & $\varepsilon_0$ ($t/U\le$) & $S$\\
\hline
$1/10$ & $1$ & $61/4000$ & $85000$ & $173/1000000\approx1/5780$ & $0.121843$\\
$1/4$ & $1$ & $51/4000$ & $102000$ & $433/1000000\approx1/2309$ & $0.121802$\\
$1/2$ & $1$ & $17/2000$ & $152000$ & $869/1000000\approx1/1151$ & $0.122001$\\
$3/4$ & $1$ & $17/4000$ & $305000$ & $1301/1000000\approx1/769$ & $0.121901$\\
$9/10$ & $1$ & $7/4000$ & $765000$ & $1537/1000000\approx1/651$ & $0.121445$\\
\hline
\end{tabular}
\end{center}
\end{theorem}

\begin{proof}
Split $B_e=B_{e,0}+B_{e,+}+B_{e,-}$ according to the change of the odd-occupation
count, $\Delta\Nodd\in\{0,2,-2\}$. The number-conserving part obeys
$B_{e,0}\le\Nodd{}_{u}+\Nodd{}_{v}$ because it can only move an electron between two
singly occupied sites or between two doubly occupied sites while the doubly occupied
count is unchanged. On the four-dimensional pair-flipping sector the two remaining
parts couple with matrix element two. Applying $2ab\le ha^{2}+b^{2}/h$ for any $h>0$
gives $g(B_{e,+}+B_{e,-})\le2gh\,\mathbf{1}+(g/h)(\Nodd{}_{u}+\Nodd{}_{v})$, and
summing over $e\in E$, each vertex having degree at most four, gives
\[
gB_E\le\Bigl(4g+\frac{4g}{h}\Bigr)\Nodd{}_{V(E)}+2kgh\,\mathbf{1}.
\]
Choosing $16g/(hU)=(1-\kappa)\gamma$, that is
$h=16g/(\gamma(1-\kappa)U)$, and using \eqref{eq:kappa}, the coefficient of
$\Nodd$ is at most $\gamma U/4$ and the scalar coefficient is
$2kgh=32kg^{2}/(\gamma(1-\kappa)U)$, which is \eqref{eq:twoab}. With
$g=\sqrt{C_{2ab}y}\,t$ and $\beta t^{2}/U\le1/C_{2ab}$ the scalar exponent per edge is
$32y/(\gamma(1-\kappa))$, and the resulting edge factor in the vertex sum is
$2Be^{32y/(\gamma(1-\kappa))}/\sqrt{C_{2ab}y}$, which is what the five certified
points test.
\end{proof}

Theorem~\ref{thm:weighted} shows that the weighted route is dominated on both counts.
The scalar exponent forces the edge factor to be large, so Koteck\'y--Preiss
convergence at the optimal split needs a scale $\ell=\sqrt{C_{2ab}y}$ of order forty,
and then \eqref{eq:kappa} caps the admissible ratio at
$\varepsilon_0=\kappa\gamma/(16\ell)$, which at the best certified point is
$1537/1000000$. Reaching $\varepsilon_0=1/50$ would need $\ell\le 45/16\approx 2.81$
at $\kappa=9/10$ and $\gamma=1$, roughly fourteen times smaller than the
vertex sum permits, and the ceiling would in addition be at least $765000$. The
comparison is thus not a matter of constants: the exact bond spectrum of
Section~\ref{sec:bond} removes the hypothesis altogether, and it also lowers the
ceiling by two orders of magnitude. The obstruction above is a structural estimate
based on the certified points, not a certified bound; what is certified is the
five-point curve and the quadratic form \eqref{eq:twoab}.

\section{What the scheme can and cannot reach}\label{sec:floors}

Theorem~\ref{thm:main} proves $6000$; the question of how much of that number is
model and how much is method is answered by optimising the continuous parameters of
the same scheme with the combinatorial tables held fixed. The first and fourth items
below are numerical and the second and third are heuristic; none of them is part of
Theorem~\ref{thm:main}.

\begin{enumerate}
\item If $B$, $q_{\mathrm{pair}}$, $T_0$ and the class scales $y_E$ are optimised
numerically with the twenty classes and the four-ary tail unchanged, the floor is
about $5314$ on all tori, at $B\approx1.112$ and $T_0\approx1.218$, and about $4000$
on even tori, at $B\approx1.114$ and $T_0\approx1.218$. The binding constraint is the
tail test \eqref{eq:tailtest}, which enters quadratically in $c_{\mathrm{tail}}$
through the argument $c_{\mathrm{tail}}\sqrt{y_tb}$ inside the exponentials of
\eqref{eq:sup}. Keeping the class scales at their certified values and optimising
only $B$ and $T_0$ gives about $5950$ on all tori, at $B\approx1.108$ and
$T_0\approx1.2115$, and about $4440$ on even tori, at $B\approx1.113$ and
$T_0\approx1.2124$, consistently with the certified failures at $C=5800$ and
$C=4300$; releasing the scales recovers the smaller numbers quoted.
\item If the coloured H\"older decoupling were replaced by a single-bond bound at
$c=1$, keeping the four-ary lattice count and the Koteck\'y--Preiss radius of about
$1.144$, the floor would drop to about $722$ (heuristic).
\item If in addition the lattice branching factor $1/w_{\mathrm{eff}}$ were replaced
by one, the floor would drop to about $9.84$ (heuristic).
\item Even for a single isolated bond, with $B=1$, no lattice sum and no colour
decoupling, matching $P_e(1,y)$ against $\ell=\sqrt{Cy}$ requires
$C\ge\inf_{y>0}P_e(1,y)^{2}/y$, which numerically is about $7.96$ on $b\in[0,40]$
(attained at $y\approx0.270$; $8.17$ at $y\approx0.258$ over all $b\ge0$); for a
single two-hop bridge, where parity supplies two factors of $\ell$ rather than one,
the corresponding floor is $\inf_{y>0}P_e(1,y)/y\approx4.78$ on $b\in[0,40]$
(attained at $y\approx0.428$; $4.97$ at $y\approx0.401$ over all $b\ge0$). Both
reported values were recomputed independently from the interval expression behind
Corollary~\ref{cor:pe}.
\end{enumerate}

The pattern is clear and it is the honest answer to the obvious question. The gap
between $6000$ and the four-ary combinatorial ideal of the third item, about a factor
of $610$, is not spent on the lattice count; it is spent on the colour amplification
inside the single-bond supremum, and item 4 isolates it: the isolated-bond floor of
$7.96$ grows like $c^{2}$ in the colour count, so at $c=5$ it is about two hundred,
and multiplied by the lattice and radius factors it accounts for the order of
magnitude. Reaching a constant of order $10$ to $100$ would require removing the
noncommutative colour decoupling, that is bounding the signless trace of several
mutually noncommuting hopping terms without the Schatten H\"older step, or weakening
the signless trajectory majorant itself. Neither is done here, and the trade-off
curve of Section~\ref{sec:tradeoff} shows that moving $q$ away from $99/100$ raises
the certified ceiling, and the numerical floor of the first item puts every
optimisation inside the scheme at $4000$ or above.

\section{Relation to the other bounds}\label{sec:relation}

Three comparisons are worth recording, each with its own caveat.

Against the earlier two-dimensional release \cite{rens2026a}: same model, same ensemble,
both unconditional in $t/U$. The ensemble is literally the same: the earlier release
writes the chemical term as $-\mu\sum_x n_x$ rather than
$-\mu\sum_x(n_x-1)$, which multiplies the partition function by the constant
$e^{\beta\mu V}$ and leaves every correlation unchanged. The ceiling improves from
$2^{27}=134\,217\,728$ to $6000$ on all tori, a factor of $22\,369$, and to $4500$ on
even tori, a factor of $29\,826$; the prefactor improves from $240/59\approx4.07$ to
$0.114$, a factor of $35.68$; the ratio $q$ rises from $5/64$ to $99/100$, so the
exponential rate $|\ln q|$ falls from $2.549$ to $0.0101$, a factor of $254$, and the
correlation length grows from $0.39$ to about a hundred lattice spacings.

Against the dimension-tracking preprint \cite{rens2026b}: the model and the ensemble
are the same, the two-dimensional constant there is $C(2)=15\,925\,248$ with
$q(2)=5/24$ and $A(2)=360/19$, so the ceiling of this paper is smaller by $2654$ on
all tori and $3539$ on even tori, the prefactor is smaller by a factor of $166$, and
the exponential rate $|\ln q|=0.0101$ is smaller by a factor of $156$ than
$\ln(24/5)=1.5686$. In three dimensions the dimension-tracking
route gives $63\,700\,992$. A two-dimensional argument of the present kind does not
transfer without redoing the colouring and the tree count; applying the same chain
formally to the cubic torus, which is not done here, would give heuristic thresholds
of about $9.2\times10^{4}$ on even cubic tori. That number is not proved, and we
therefore make no claim in three dimensions.

Against the physical superexchange scale: the hopping-induced exchange scale that
governs pair formation in the strong-coupling regime is of order $t^{2}/U$
per bond, so a coordination-based estimate of the clustering
onset is a few tens of $t^{2}/U$ on the square lattice. The ceiling $6000\,t^{2}/U$ is
a rigorous upper bound and is about two orders of magnitude above that estimate. This
comparison is heuristic in both directions: the superexchange scale is not a
clustering temperature and the ceiling is not a physical transition temperature, the
gap between them quantifies the remaining slack of the method rather than a
discrepancy in the physics, and we draw no conclusion about where pair correlations
actually decay.

\section{Discussion and limitations}\label{sec:discussion}

What is proved is the following. At temperatures above $6000\,t^{2}/U$ on any square
torus, and above $4500\,t^{2}/U$ on tori of even side, equal-time on-site pair
correlations decay exponentially in the torus distance with an explicit prefactor and
an explicit rate, uniformly in the phonon frequency, in the torus side and in the
chemical potential; the rows are uniformly summable, and the statement holds in
canonical thermodynamic limits at interior density. The argument imposes no smallness
condition on $t/U$ and does not assume the adiabatic limit.

The limitations are equally explicit and each is a property of what is written here,
not a conjecture.

The constant is a bound on a clustering temperature of one correlator. It is not the
critical temperature of anything, no phase transition is proved or excluded, and
nothing here shows that the model is superconducting at any temperature. The decay
rate is weak, $|\ln q|=0.01005$, so the bound should be read as a statement about
uniform summability of the correlation matrix and about an exact threshold in
$\beta t^{2}/U$, not as a quantitative decay estimate at moderate distance.

The proof covers the two-dimensional square lattice. The colouring, the small-edge
enumeration and the four-ary tree count are all specific to the square grid; the
dimension-tracking paper \cite{rens2026b} treats the general cubic case by a different
route with a much larger constant and a better decay ratio, and no dimension-uniform
improvement follows from the present argument.

The constants $6000$ and $4500$ are not optimal, and Section~\ref{sec:floors}
quantifies how far the method can be pushed without a new idea: to about $5314$ and
$4000$ by parameter optimisation, and no further while the coloured H\"older
decoupling and the signless trajectory majorant are both retained. The certified
failures at $C=5800$ on all tori and at $C=4300$ on even tori are failures of the
particular tail test, and they are evidence about the scheme, not about the model.

Finally, the certificate program proves arithmetic, not analysis. Its output asserts
exact rational and symbolic identities and inequalities, including the interval
bounds of Corollary~\ref{cor:pe} and the budget comparisons of
Table~\ref{tab:audit}; it does not replace the operator inequalities of
Lemmas~\ref{lem:bridge}, \ref{lem:holder} and \ref{lem:counting} or the estimates of
Theorems~\ref{thm:bond} and~\ref{thm:weighted}, each of which has a written
derivation in this paper.

\section{Supporting material}\label{sec:supplement}

The supplement is a single program of about eight hundred lines that recomputes, in
exact rational and symbolic arithmetic, every finite witness quoted above; it needs
no data file. Run without arguments it prints seven certified lines and exits
zero only if all of them hold: the sixteen-state bond spectrum, the multiplicities
and the two half-filling identities of Theorem~\ref{thm:bond}; the exact torus
enumeration at $k\le6$, the four-ary coefficients \eqref{eq:raney} and the
two-variable tree check; the $102$ certified intervals of Corollary~\ref{cor:pe}
over the twenty-two class slots; the two budgets and the improvement factors of
Theorem~\ref{thm:main}; and the ten points of the trade-off curve. It also verifies
the five certified points of Theorem~\ref{thm:weighted}. Run with the flag
\verb|--mutants| it repeats the same checks with one guarded claim corrupted at a
time, twelve times, and exits zero only if every corrupted claim is rejected. Two of
the corrupted claims are the ceilings $C=5800$ on all tori and $C=4300$ on even
tori, which fail the tail test \eqref{eq:tailtest}; the others corrupt the bridge
parity, a colour count, the tail colour on odd tori, one interval bound, the
half-filling monotonicity, the even-torus class table, the canonical source weight,
the pair prefactor, the weighted obstruction and the tree coefficients.

The certified run and the mutant run were carried out on a clean copy of the
supplement with uv 0.9.5, Python 3.13.9, sympy 1.14.0 and mpmath 1.3.0, and both
exited zero. The manifest in the supplement lists the SHA-256 digests of its files;
the certificate program has digest
\texttt{3449ce2455935b29c03616cb3c3ed70018903d9ce32d59f0d1b27f93dfe7f62f}.

\section*{AI assistance}

The analytical results, the design of the estimates, and the final text are the
author's responsibility. Independent proof reviews of the argument and of the
certificate program were carried out by separate AI systems, whose comments led to
several corrections recorded in the present version. The manuscript, its LaTeX
source, and the certificate program were drafted and edited with AI assistance under
the author's direction and review.
\addcontentsline{toc}{section}{AI assistance}

This preprint has not been peer reviewed.

\appendix

\section{The residual oscillator weight}\label{app:oscillator}

Section~\ref{sec:langfirsov} uses the exact reduction of the oscillator trace to the
contraction \eqref{eq:majorant}. This appendix repeats the statement in the units and
notation of this paper, records the kernel that enters it, and derives the mode form
\eqref{eq:phi}. Only positivity of the weight is used in the proof of
Theorem~\ref{thm:main}, and no step below constrains $\Omega$.

Throughout $\hbar=m=1$, so that $K=\Omega^{2}$, $U=F^{2}/\Omega^{2}$ and
$\alpha=\alpha_{\mathrm{LF}}=\sqrt{U/(2\Omega)}$; the site label is dropped and
$H_{\mathrm{ph}}=\Omega(b^{\dagger}b+\tfrac12)$. Let $n(u)\in\{0,1,2\}$ be the
occupation of the site along a closed trajectory, that is a $\beta$-periodic,
right continuous, piecewise constant function of the imaginary time
$u\in[0,\beta)$ with jumps $c_j=n(u_j^{+})-n(u_j^{-})\in\{\pm1,\pm2\}$ at
$0\le u_1<\dots<u_m<\beta$ and $\sum_jc_j=0$. Write
$D^{c}(u)=\exp\bigl(c\alpha(e^{\Omega u}b^{\dagger}-e^{-\Omega u}b)\bigr)$ for the
analytically continued displacement carried by a hop, and let
$\langle\cdot\rangle_{\mathrm{ph}}$ be the trace over the oscillator normalized by
$\mathrm{Tr}\,e^{-\beta H_{\mathrm{ph}}}$.

\begin{lemma}[chronological displacement trace]\label{lem:chronological}
Let
\begin{equation}\label{eq:kernel}
K_\beta(w)=\frac{U\Omega}{2}\,\frac{e^{-\Omega w}+e^{-\Omega(\beta-w)}}{1-e^{-\beta\Omega}}
=\frac{U\Omega}{2}\,\frac{\cosh\bigl(\Omega(\tfrac{\beta}{2}-w)\bigr)}{\sinh(\beta\Omega/2)},
\qquad 0\le w\le\beta,
\end{equation}
extended $\beta$-periodically to $w\in\mathbb{R}$, and
\begin{equation}\label{eq:phip}
\Phi_\beta[n]=\frac14\int_0^\beta\!\!\int_0^\beta\bigl(n(u)-n(v)\bigr)^2
K_\beta(u-v)\,du\,dv .
\end{equation}
Then, the factors of the time-ordered product being ordered from right to left in
increasing time,
\begin{equation}\label{eq:chronological}
\Bigl\langle\mathcal{T}\prod_{j}D^{c_j}(u_j)\Bigr\rangle_{\mathrm{ph}}
=\exp\bigl(-\Phi_\beta[n]\bigr)
=\exp\Bigl(\frac12\iint n(u)n(v)K_\beta(u-v)\,du\,dv-\frac{U}{2}\int_0^\beta n(u)^2du\Bigr).
\end{equation}
The kernel is positive and $\int_0^\beta K_\beta(w)\,dw=U$; consequently
$0\le\exp(-\Phi_\beta[n])\le1$, with equality if and only if $n$ is constant.
Finally, with the modes $n_s=\beta^{-1}\int_0^\beta n(u)e^{-i\omega_su}du$ and
$\omega_s=2\pi s/\beta$,
\begin{equation}\label{eq:phimode}
\Phi_\beta[n]=\frac{\beta U}{2}\sum_{s\ne0}
\frac{\omega_s^{2}}{\omega_s^{2}+\Omega^{2}}\,|n_s|^{2}.
\end{equation}
\end{lemma}

\begin{proof}
The identity \eqref{eq:chronological} is proved in \cite[Section 3]{rens2026a}: the
displacement factors are normal ordered with the Baker--Campbell--Hausdorff formula,
the thermal occupations $\bar n=(e^{\beta\Omega}-1)^{-1}$ are summed as geometric
series, and the kernel \eqref{eq:kernel} is the covariance of the oscillator at
imaginary-time separation $w$. The two exponents in \eqref{eq:chronological}
coincide by the mass identity derived below.

For \eqref{eq:phimode} compute the Fourier coefficients
$K_s=\int_0^\beta K_\beta(w)e^{-i\omega_sw}dw$ of the kernel on the circle of
circumference $\beta$, so that $K_0$ is its mass. Integrating
$e^{-(\Omega\pm i\omega_s)w}$ over $[0,\beta]$ gives
\begin{equation}\label{eq:kernelmode}
K_0=\frac{U\Omega}{2}\Bigl(\frac1\Omega+\frac1\Omega\Bigr)=U,
\qquad
K_s=\frac{U\Omega^{2}}{\Omega^{2}+\omega_s^{2}}\quad(s\ne0).
\end{equation}
Expand $n(u)=\sum_sn_se^{i\omega_su}$. Then $\int_0^\beta n^{2}du
=\beta\sum_s|n_s|^{2}$ and $\iint n(u)n(v)K_\beta(u-v)\,dudv
=\beta\sum_s|n_s|^{2}K_s$, so
\begin{equation}
\frac14\iint(n-n')^{2}K_\beta
=\frac{U}{2}\int_0^\beta n^{2}du-\frac12\iint nK_\beta n
=\frac{\beta}{2}\sum_s|n_s|^{2}\bigl(U-K_s\bigr)
=\frac{\beta U}{2}\sum_{s\ne0}\frac{\omega_s^{2}}{\omega_s^{2}+\Omega^{2}}|n_s|^{2},
\end{equation}
where the first equality is the mass identity and the last uses $U-K_0=0$. By
\eqref{eq:kernel} the kernel is a positive constant times a positive hyperbolic
cosine, so the integral \eqref{eq:phip} is nonnegative and vanishes only if $n$ is
constant, and \eqref{eq:kernelmode} is the mass identity $\int_0^\beta K_\beta=U$.
\end{proof}

Two remarks on the use made of the lemma. First, the weight is local in the site
index, because the oscillator at $x$ is driven only by $n_x$; the product of the
weights \eqref{eq:chronological} over the sites of a closed electronic trajectory is
its phonon weight, and \eqref{eq:majorant} holds on any lattice, in particular on the
square torus. Second, \eqref{eq:chronological} exhibits the weight as a product of the
nonnegative oscillator functional $\exp\bigl(\tfrac12\iint nK_\beta n\bigr)$ and of
the atomic attraction $\exp\bigl(-\tfrac U2\int n^{2}\bigr)$, and it is the positivity
of the kernel in the dipole form \eqref{eq:phip}, not any smallness in $\Omega$ or in
$\beta$, that makes the product at most one.

\bibliographystyle{plain}
\bibliography{refs}
\end{document}
